% GENERATED FILE. Edit canonical lessons, then run npm run generate:books.
% canonical-source-hash: fnv1a64:39f65c0127b434ad
% canonical-lessons: ML-C508-caruka-drizzle, ML-C508-pemari, ML-C508-nanayuka, ML-C508-malakkottu, ML-C508-unakkuka, ML-R508-review-501-508

\chapter{Getting Wet}
\label{ch:ml-getting-wet}
\hlchaptermodality{508}

\begin{chapteropening}
\textbf{By the end of this chapter:} \emph{I can say to drizzle, a downpour, to get wet, a raincoat and to dry out.}

Complete the last lesson of chapter 508: I can say to drizzle, a downpour, to get wet, a raincoat and to dry out.
\end{chapteropening}

\section[\texorpdfstring{\ml{ചാറുക} (cāṟuka)}{ചാറുക (cāṟuka)}]{\ml{ചാറുക} (cāṟuka) — to drizzle}

\label{lesson:ML-C508-caruka-drizzle}



\begin{quote}
Before the new one: say the Malayalam for to stop raining, then the Malayalam for a light drizzle.
\end{quote}

\subsection*{You'll want to know: \ml{ചാറുക}}
\textbf{\ml{ചാറുക}} — \emph{cāṟuka} — \textquotedblleft{}to drizzle\textquotedblright{}.

\begin{quote}
\ml{പുറത്ത്} \ml{മഴ} \ml{ചാറുന്നു}.
\end{quote}

\begin{quote}
\emph{puṟattŭ maḻa cāṟunnu.}
\end{quote}

\begin{quote}
It is drizzling outside.
\end{quote}

Past: \ml{ചാറി}, \emph{cāṟi}.

The subject is the rain: \ml{മഴ} \ml{ചാറുന്നു}, the rain drizzles.

\begin{cousinweb}
\textbf{\ml{ചാറുക}} (\emph{cāṟuka}) is native Dravidian, \textquotedblleft{}to run off; to drizzle\textquotedblright{} (DEDR 2482).
\end{cousinweb}

\begin{grammarlens}[title={What you've built}]
One more word for getting around.
\end{grammarlens}

\subsection*{Your turn}
\begin{itemize}
\raggedright
  \item \emph{Say it:} \emph{cāṟuka}
  \item \emph{Say it:} \emph{cāṟuka}, once more
  \item \emph{Say it:} \emph{puṟattŭ maḻa cāṟunnu.} — It is drizzling outside.
  \item \emph{Recall:} say the Malayalam for to stop raining, then the Malayalam for a light drizzle, then say \emph{cāṟuka} again
\end{itemize}

\begin{tcolorbox}[breakable,colback=teal!4,colframe=teal!35!black,title={Before you move on}]
What does \ml{ചാറുക} mean? (\textquotedblleft{}To drizzle\textquotedblright{}.) Say it once more.
\end{tcolorbox}
\section[\texorpdfstring{\ml{പേമാരി} (pēmāri)}{പേമാരി (pēmāri)}]{\ml{പേമാരി} (pēmāri) — a downpour; torrential rain}

\label{lesson:ML-C508-pemari}



\begin{quote}
Before the new one: say the Malayalam for a light drizzle, then the Malayalam for to drizzle.
\end{quote}

\subsection*{You'll want to know: \ml{പേമാരി}}
\textbf{\ml{പേമാരി}} — \emph{pēmāri} — \textquotedblleft{}a downpour; torrential rain\textquotedblright{}.

\begin{quote}
\ml{രാത്രി} \ml{മുഴുവൻ} \ml{പേമാരി} \ml{ആയിരുന്നു}.
\end{quote}

\begin{quote}
\emph{rātri muḻuvan pēmāri āyirunnu.}
\end{quote}

\begin{quote}
There was torrential rain all night.
\end{quote}

Stronger than \ml{മഴ}: the rain of a storm.

\begin{cousinweb}
\textbf{\ml{പേമാരി}} (\emph{pēmāri}) is a Malayalam word for heavy rain whose second part is \ml{മാരി}, heavy rain (DEDR 4819, which compares Sanskrit \emph{māri}, rain); the origin of the first part is not settled.
\end{cousinweb}

\begin{grammarlens}[title={What you've built}]
One more word for getting around.
\end{grammarlens}

\subsection*{Your turn}
\begin{itemize}
\raggedright
  \item \emph{Say it:} \emph{pēmāri}
  \item \emph{Say it:} \emph{pēmāri}, once more
  \item \emph{Say it:} \emph{rātri muḻuvan pēmāri āyirunnu.} — There was torrential rain all night.
  \item \emph{Recall:} say the Malayalam for a light drizzle, then the Malayalam for to drizzle, then say \emph{pēmāri} again
\end{itemize}

\begin{tcolorbox}[breakable,colback=teal!4,colframe=teal!35!black,title={Before you move on}]
What does \ml{പേമാരി} mean? (\textquotedblleft{}A downpour; torrential rain\textquotedblright{}.) Say it once more.
\end{tcolorbox}
\section[\texorpdfstring{\ml{നനയുക} (nanayuka)}{നനയുക (nanayuka)}]{\ml{നനയുക} (nanayuka) — to get wet}

\label{lesson:ML-C508-nanayuka}



\begin{quote}
Before the new one: say the Malayalam for to drizzle, then the Malayalam for a downpour.
\end{quote}

\subsection*{You'll want to know: \ml{നനയുക}}
\textbf{\ml{നനയുക}} — \emph{nanayuka} — \textquotedblleft{}to get wet\textquotedblright{}.

\begin{quote}
\ml{കുട} \ml{ഇല്ലാത്തതുകൊണ്ട്} \ml{ഞാൻ} \ml{നനഞ്ഞു}.
\end{quote}

\begin{quote}
\emph{kuṭa illāttatukoṇṭŭ ñān nanaññu.}
\end{quote}

\begin{quote}
Because I had no umbrella, I got wet.
\end{quote}

Past: \ml{നനഞ്ഞു}, \emph{nanaññu}.

The person who gets wet is the subject: \ml{ഞാൻ} \ml{നനഞ്ഞു}, I got wet.

\begin{cousinweb}
\textbf{\ml{നനയുക}} (\emph{nanayuka}) is native Dravidian, \textquotedblleft{}to become wet, to be soaked\textquotedblright{} (DEDR 3630).
\end{cousinweb}

\begin{grammarlens}[title={What you've built}]
One more word for getting around.
\end{grammarlens}

\subsection*{Your turn}
\begin{itemize}
\raggedright
  \item \emph{Say it:} \emph{nanayuka}
  \item \emph{Say it:} \emph{nanayuka}, once more
  \item \emph{Say it:} \emph{kuṭa illāttatukoṇṭŭ ñān nanaññu.} — Because I had no umbrella, I got wet.
  \item \emph{Recall:} say the Malayalam for to drizzle, then the Malayalam for a downpour, then say \emph{nanayuka} again
\end{itemize}

\begin{tcolorbox}[breakable,colback=teal!4,colframe=teal!35!black,title={Before you move on}]
What does \ml{നനയുക} mean? (\textquotedblleft{}To get wet\textquotedblright{}.) Say it once more.
\end{tcolorbox}
\section[\texorpdfstring{\ml{മഴക്കോട്ട്} (maḻakkōṭṭŭ)}{മഴക്കോട്ട് (maḻakkōṭṭŭ)}]{\ml{മഴക്കോട്ട്} (maḻakkōṭṭŭ) — a raincoat}

\label{lesson:ML-C508-malakkottu}



\begin{quote}
Before the new one: say the Malayalam for a downpour, then the Malayalam for to get wet.
\end{quote}

\subsection*{You'll want to know: \ml{മഴക്കോട്ട്}}
\textbf{\ml{മഴക്കോട്ട്}} — \emph{maḻakkōṭṭŭ} — \textquotedblleft{}a raincoat\textquotedblright{}.

\begin{quote}
\ml{അവൻ} \ml{മഴക്കോട്ട്} \ml{ഇട്ട്} \ml{ബൈക്കിൽ} \ml{പോയി}.
\end{quote}

\begin{quote}
\emph{avan maḻakkōṭṭŭ iṭṭŭ baikkil pōyi.}
\end{quote}

\begin{quote}
He put on a raincoat and went on his motorbike.
\end{quote}

\ml{മഴ}, rain, plus English coat.

\begin{cousinweb}
\textbf{\ml{മഴക്കോട്ട്}} (\emph{maḻakkōṭṭŭ}) is a Malayalam compound built from the native \ml{മഴ}, rain (DEDR 4753), and \ml{കോട്ട്}, which was borrowed from English \emph{coat}.
\end{cousinweb}

\begin{grammarlens}[title={What you've built}]
One more word for getting around.
\end{grammarlens}

\subsection*{Your turn}
\begin{itemize}
\raggedright
  \item \emph{Say it:} \emph{maḻakkōṭṭŭ}
  \item \emph{Say it:} \emph{maḻakkōṭṭŭ}, once more
  \item \emph{Say it:} \emph{avan maḻakkōṭṭŭ iṭṭŭ baikkil pōyi.} — He put on a raincoat and went on his motorbike.
  \item \emph{Recall:} say the Malayalam for a downpour, then the Malayalam for to get wet, then say \emph{maḻakkōṭṭŭ} again
\end{itemize}

\begin{tcolorbox}[breakable,colback=teal!4,colframe=teal!35!black,title={Before you move on}]
What does \ml{മഴക്കോട്ട്} mean? (\textquotedblleft{}A raincoat\textquotedblright{}.) Say it once more.
\end{tcolorbox}
\section[\texorpdfstring{\ml{ഉണക്കുക} (uṇakkuka)}{ഉണക്കുക (uṇakkuka)}]{\ml{ഉണക്കുക} (uṇakkuka) — to dry out; to hang out to dry}

\label{lesson:ML-C508-unakkuka}



\begin{quote}
Before the new one: say the Malayalam for to get wet, then the Malayalam for a raincoat.
\end{quote}

\subsection*{You'll want to know: \ml{ഉണക്കുക}}
\textbf{\ml{ഉണക്കുക}} — \emph{uṇakkuka} — \textquotedblleft{}to dry out; to hang out to dry\textquotedblright{}.

\begin{quote}
\ml{വെയിൽ} \ml{വന്നപ്പോൾ} \ml{അമ്മ} \ml{തുണികൾ} \ml{ഉണക്കി}.
\end{quote}

\begin{quote}
\emph{veyil vannappōḷ amma tuṇikaḷ uṇakki.}
\end{quote}

\begin{quote}
When the sun came out, mother dried the clothes.
\end{quote}

Past: \ml{ഉണക്കി}, \emph{uṇakki}.

The transitive partner of \ml{ഉണങ്ങുക}, to become dry.

\begin{cousinweb}
\textbf{\ml{ഉണക്കുക}} (\emph{uṇakkuka}) is native Dravidian, \textquotedblleft{}to put to dry, to air\textquotedblright{} (DEDR 601).
\end{cousinweb}

\begin{grammarlens}[title={What you've built}]
One more word for getting around.
\end{grammarlens}

\subsection*{Your turn}
\begin{itemize}
\raggedright
  \item \emph{Say it:} \emph{uṇakkuka}
  \item \emph{Say it:} \emph{uṇakkuka}, once more
  \item \emph{Say it:} \emph{veyil vannappōḷ amma tuṇikaḷ uṇakki.} — When the sun came out, mother dried the clothes.
  \item \emph{Recall:} say the Malayalam for to get wet, then the Malayalam for a raincoat, then say \emph{uṇakkuka} again
\end{itemize}

\begin{tcolorbox}[breakable,colback=teal!4,colframe=teal!35!black,title={Before you move on}]
What does \ml{ഉണക്കുക} mean? (\textquotedblleft{}To dry out; to hang out to dry\textquotedblright{}.) Say it once more.
\end{tcolorbox}
\section[(dialogue)]{Words from chapters 501–508 — a review}

\label{lesson:ML-R508-review-501-508}



\begin{quote}
Say the Malayalam for a zoo and a public garden.
\end{quote}

\subsection*{Your turn}
Say these aloud:

\begin{itemize}
\raggedright
  \item a zoo, a public garden, a whale, an entry fee, to drift in
  \item a plantation, greenery, to sprout, to fall off, to cut down
  \item the environment, surroundings, pollution, to pollute, sewage
  \item litter, to throw away, to pile up, to be ruined, to reuse
  \item energy, a shortage, groundwater, global warming, a rain pit
  \item sunrise, sunset, the horizon, the atmosphere, rain clouds
  \item the monsoon, the southwest monsoon, the northeast monsoon, to stop raining, a light drizzle
  \item to drizzle, a downpour, to get wet, a raincoat, to dry out
\end{itemize}

\begin{tcolorbox}[breakable,colback=teal!4,colframe=teal!35!black,title={Before you move on}]
A zoo? (\textbf{\ml{മൃഗശാല}}, \emph{mṛgaśāla}.) A public garden? (\textbf{\ml{ഉദ്യാനം}}, \emph{udyānaṁ}.) A whale? (\textbf{\ml{തിമിംഗലം}}, \emph{timiṁgalaṁ}.) An entry fee? (\textbf{\ml{പ്രവേശനഫീസ്}}, \emph{pravēśanaphīsŭ}.) To drift in? (\textbf{\ml{ഒഴുകിയെത്തുക}}, \emph{oḻukiyettuka}.) A plantation? (\textbf{\ml{തോട്ടം}}, \emph{tōṭṭaṁ}.) Greenery? (\textbf{\ml{പച്ചപ്പ്}}, \emph{paccappŭ}.) To sprout? (\textbf{\ml{തളിർക്കുക}}, \emph{taḷirkkuka}.) To fall off? (\textbf{\ml{കൊഴിയുക}}, \emph{koḻiyuka}.) To cut down? (\textbf{\ml{വെട്ടിമാറ്റുക}}, \emph{veṭṭimāṟṟuka}.) The environment? (\textbf{\ml{പരിസ്ഥിതി}}, \emph{paristhiti}.) Surroundings? (\textbf{\ml{പരിസരം}}, \emph{parisaraṁ}.) Pollution? (\textbf{\ml{മലിനീകരണം}}, \emph{malinīkaraṇaṁ}.) To pollute? (\textbf{\ml{മലിനമാക്കുക}}, \emph{malinamākkuka}.) Sewage? (\textbf{\ml{മലിനജലം}}, \emph{malinajalaṁ}.) Litter? (\textbf{\ml{ചവറ്}}, \emph{cavaṟŭ}.) To throw away? (\textbf{\ml{വലിച്ചെറിയുക}}, \emph{valicceṟiyuka}.) To pile up? (\textbf{\ml{കുമിഞ്ഞുകൂടുക}}, \emph{kumiññukūṭuka}.) To be ruined? (\textbf{\ml{നശിക്കുക}}, \emph{naśikkuka}.) To reuse? (\textbf{\ml{പുനരുപയോഗിക്കുക}}, \emph{punarupayōgikkuka}.) Energy? (\textbf{\ml{ഊർജ്ജം}}, \emph{ūrjjaṁ}.) A shortage? (\textbf{\ml{ക്ഷാമം}}, \emph{kṣāmaṁ}.) Groundwater? (\textbf{\ml{ഭൂഗർഭജലം}}, \emph{bhūgarbhajalaṁ}.) Global warming? (\textbf{\ml{ആഗോളതാപനം}}, \emph{āgōḷatāpanaṁ}.) A rain pit? (\textbf{\ml{മഴക്കുഴി}}, \emph{maḻakkuḻi}.) Sunrise? (\textbf{\ml{സൂര്യോദയം}}, \emph{sūryōdayaṁ}.) Sunset? (\textbf{\ml{സൂര്യാസ്തമയം}}, \emph{sūryāstamayaṁ}.) The horizon? (\textbf{\ml{ചക്രവാളം}}, \emph{cakravāḷaṁ}.) The atmosphere? (\textbf{\ml{അന്തരീക്ഷം}}, \emph{antarīkṣaṁ}.) Rain clouds? (\textbf{\ml{മഴക്കാർ}}, \emph{maḻakkār}.) The monsoon? (\textbf{\ml{മൺസൂൺ}}, \emph{maṇsūṇ}.) The southwest monsoon? (\textbf{\ml{കാലവർഷം}}, \emph{kālavarṣaṁ}.) The northeast monsoon? (\textbf{\ml{തുലാവർഷം}}, \emph{tulāvarṣaṁ}.) To stop raining? (\textbf{\ml{തോരുക}}, \emph{tōruka}.) A light drizzle? (\textbf{\ml{ചാറ്റൽമഴ}}, \emph{cāṟṟalmaḻa}.) To drizzle? (\textbf{\ml{ചാറുക}}, \emph{cāṟuka}.) A downpour? (\textbf{\ml{പേമാരി}}, \emph{pēmāri}.) To get wet? (\textbf{\ml{നനയുക}}, \emph{nanayuka}.) A raincoat? (\textbf{\ml{മഴക്കോട്ട്}}, \emph{maḻakkōṭṭŭ}.) To dry out? (\textbf{\ml{ഉണക്കുക}}, \emph{uṇakkuka}.)
\end{tcolorbox}
